% ===== PRÉAMBULE CANONIQUE — à coller VERBATIM en tête de CHAQUE .tex =====
\documentclass[11pt,a4paper]{article}
\usepackage[utf8]{inputenc}\usepackage[T1]{fontenc}\usepackage{lmodern}
\usepackage[french]{babel}
\usepackage{amsmath,amssymb,mathtools}\usepackage{icomma}
\usepackage{siunitx}\sisetup{locale=FR}  % \qty{}{}, \si{}, \ang{}
\usepackage{xcolor}\usepackage{colortbl}
\usepackage{tikz}\usepackage{pgfplots}\pgfplotsset{compat=1.18}
\usepackage[siunitx,europeanresistors]{circuitikz} % circuits (élec.)
\usepackage[most]{tcolorbox}\usepackage{enumitem}\usepackage{array,booktabs}
\usepackage[a4paper,margin=2.2cm]{geometry}
\usetikzlibrary{arrows.meta,calc,angles,quotes,positioning,patterns,%
  backgrounds,shapes.geometric,decorations.pathmorphing,decorations.markings,intersections}
\definecolor{primary}{RGB}{222,49,99}\definecolor{primarydark}{RGB}{150,22,55}
\definecolor{defcol}{RGB}{0,114,178}\definecolor{methcol}{RGB}{52,92,148}
\definecolor{excol}{RGB}{0,135,90}\definecolor{attncol}{RGB}{200,40,55}
\tcbset{boxbase/.style={enhanced,breakable,boxrule=0.9pt,arc=2.5pt,
  left=3mm,right=3mm,top=2.3mm,bottom=2.3mm,fonttitle=\bfseries,coltitle=white}}
% --- boîtes TUTORIEL (noms EXACTS, ne pas renommer) ---
\newtcolorbox{cle}[1][]{boxbase,colback=defcol!6,colframe=defcol,
  title={\textsc{À retenir}\ifx\relax#1\relax\relax\else\textmd{ : \itshape #1}\fi}}
\newtcolorbox{methode}[1][]{boxbase,colback=methcol!7,colframe=methcol,sharp corners,
  title={\textsc{Méthode}\ifx\relax#1\relax\relax\else\textmd{ --- \itshape #1}\fi}}
\newtcolorbox{exemplebox}[1][]{boxbase,colback=excol!5,colframe=excol,
  title={\textsc{Exemple}\ifx\relax#1\relax\relax\else\textmd{ : \itshape #1}\fi}}
\newtcolorbox{piege}[1][]{boxbase,colback=attncol!6,colframe=attncol,
  title={\textsc{Piège}\ifx\relax#1\relax\relax\else\textmd{ : \itshape #1}\fi}}
\newtcolorbox{remarque}[1][]{boxbase,colback=black!4,colframe=black!45,
  title={\textsc{Remarque}\ifx\relax#1\relax\relax\else\textmd{ : \itshape #1}\fi}}
% --- boîtes EXERCICE / CORRIGÉ (noms EXACTS, auto-comptées) ---
\newtcolorbox[auto counter]{exo}[1][]{boxbase,colback=primary!4,colframe=primary,
  title={\textsc{Exercice}~\thetcbcounter\ifx\relax#1\relax\relax\else\textmd{ --- \itshape #1}\fi}}
\newtcolorbox[auto counter]{corrige}[1][]{boxbase,colback=excol!4,colframe=excol,
  title={\textsc{Corrigé}~\thetcbcounter\ifx\relax#1\relax\relax\else\textmd{ --- \itshape #1}\fi}}
\newcommand{\vv}[1]{\vec{#1}}\newcommand{\ds}{\displaystyle}
\newcommand{\R}{\mathbb{R}}\newcommand{\N}{\mathbb{N}}\newcommand{\Z}{\mathbb{Z}}
% (un \usepackage{fancyhdr}+en-tête est possible ; il est ignoré côté web)
% =========================================================================
\begin{document}
\sisetup{per-mode=symbol}

\begin{corrige}[frottement : qui gagne, qui perd ?]
\begin{enumerate}[label=\alph*)]
  \item La tige est devenue \emph{négative} : elle a \textbf{gagné} des électrons. Le transfert va
        donc de la \textbf{laine vers la tige}. Nombre transféré :
        \[ N = \frac{|q_{\text{tige}}|}{e} = \frac{\num{3.2e-9}}{\num{1.6e-19}} = \num{2e10}\ \text{électrons}. \]
  \item Conservation de la charge : la laine a \emph{perdu} exactement ces $\num{2e10}$ électrons,
        elle porte donc la charge opposée
        \[ q_{\text{laine}} = +\,\num{2e10}\times e = \SI{+3.2e-9}{\coulomb}. \]
        (La somme des charges reste nulle : $q_{\text{tige}}+q_{\text{laine}}=0$.)
\end{enumerate}
\end{corrige}

\begin{corrige}[partages successifs par contact]
Deux sphères \emph{identiques} en contact se partagent également la charge totale.
\begin{enumerate}[label=\alph*)]
  \item $A$ et $B$ : charge totale $\SI{12}{\nano\coulomb}$ répartie en deux parts égales.
        \[ q_A' = q_B' = \frac{\SI{12}{\nano\coulomb}}{2} = \SI{6}{\nano\coulomb}\ \text{chacune}. \]
  \item $B$ ($\SI{6}{\nano\coulomb}$) touche $C$ (neutre) : total $\SI{6}{\nano\coulomb}$ partagé,
        \[ q_C = \frac{\SI{6}{\nano\coulomb}}{2} = \SI{3}{\nano\coulomb}. \]
\end{enumerate}
\end{corrige}

\begin{corrige}[influence puis mise à la terre]
\begin{enumerate}[label=\alph*)]
  \item Le bâton négatif \emph{repousse} les électrons libres de la sphère vers le côté éloigné.
        En reliant à la terre, ces électrons excédentaires \textbf{s'écoulent vers le sol}. On coupe
        la terre \emph{avant} de retirer le bâton : la sphère se retrouve alors avec un
        \emph{déficit} d'électrons. Signe final : \textbf{positif} (opposé au bâton).
  \item Avec un bâton \emph{positif}, il attirerait des électrons de la terre vers la sphère : signe
        final \textbf{négatif}. \emph{Règle : l'électrisation par influence avec mise à la terre
        donne une charge de signe opposé à celui de l'inducteur.}
\end{enumerate}
\end{corrige}

\begin{corrige}[charge et masse : combien d'électrons partis ?]
\begin{enumerate}[label=\alph*)]
  \item $q>0$ : la bille a \emph{perdu} des électrons.
        \[ N = \frac{|q|}{e} = \frac{\num{1e-6}}{\num{1.6e-19}} = \num{6.25e12}\ \text{électrons}. \]
  \item Chaque électron parti emporte sa masse $m_e$ ; la masse \textbf{diminue} de
        \[ \Delta m = -\,N\,m_e = -\,\num{6.25e12}\times\num{9.1e-31} \approx \SI{-5.7e-18}{\kilogram}. \]
        Variation infime : la charge modifie la masse de façon totalement négligeable.
\end{enumerate}
\end{corrige}

\begin{corrige}[bilan protons--électrons]
\begin{enumerate}[label=\alph*)]
  \item Neutre $\Rightarrow$ autant d'électrons que de protons : $N_e = N_p = \num{4.0e14}$ électrons.
  \item Il capte $\num{5.0e10}$ électrons : il devient négatif, avec un excès de $\num{5.0e10}$
        électrons.
        \[ q = -\,\num{5.0e10}\times e = -\,\num{5.0e10}\times\num{1.6e-19} = \SI{-8.0e-9}{\coulomb} = \SI{-8}{\nano\coulomb}. \]
\end{enumerate}
\end{corrige}

\end{document}
